\documentclass[11pt]{article}
\usepackage[margin=1in]{geometry}
\usepackage{amsmath,amssymb,amsthm,booktabs}
\usepackage[hidelinks]{hyperref}
\usepackage{xurl}
\setlength{\emergencystretch}{2em}
\newtheorem{lemma}{Lemma}
\newtheorem{theorem}[lemma]{Theorem}
\newtheorem{corollary}[lemma]{Corollary}
\newcommand{\F}{\mathbb F}
\newcommand{\mr}{\operatorname{mr}}
\newcommand{\rk}{\operatorname{rank}}
\title{Twenty-four sharp meridional bounds from finite affine images over $\F_3$}
\author{Written proof and exact finite certificates}
\date{17 November 2026}
\begin{document}
\maketitle

Twenty-four knot groups have explicit finite linear images whose
meridian conjugacy classes require four generators. Their diagrams
have four Wirtinger seeds, so all twenty-four have meridional rank and
bridge number four. The image groups themselves require only two
ordinary generators. The proof uses literal matrices and finite
group calculations; no branched-cover group identification or
homology inequality is a premise.

\section{Fixed spaces and meridian conjugacy classes}

For a representation of a group $H$ on $V$, write $V^H$ for the
common fixed space. The following elementary observation supplies
the lower bound.

\begin{lemma}\label{lem:fixed}
If $H$ is generated by $k$ elements whose displacement matrices
have rank at most $\delta$, then
\[
                   \dim V-\dim V^H\le k\delta.
\]
In particular, if a knot group has an image with common fixed-space
codimension $q$ and meridian displacement rank $\delta>0$, then
$\mr(K)\ge\lceil q/\delta\rceil$.
\end{lemma}
\begin{proof}
The fixed space is the intersection of the kernels of the $k$
displacement maps. Its codimension is at most the sum of their
ranks. Conjugate meridians have equal displacement rank.
\end{proof}

The lemma does not require a preserved form, a faithful action,
or an involutory meridian. Its use below should not be interpreted
as a new general linear-algebra bound.

\section{The twenty-four finite certificates}

Let
\[
\begin{split}
\mathcal K=\{&13a_{304},13a_{362},13a_{623},13a_{629},
13a_{3232},13a_{3259},13a_{3374},13a_{3375},\\
&13n_{834},13n_{838},13n_{839},13n_{991},13n_{992},13n_{993},\\
&13n_{1367},13n_{1370},13n_{1371},13n_{1374},13n_{1375},13n_{1376},\\
&13n_{3814},13n_{3840},13n_{3911},13n_{3914}\}.
\end{split}
\]
Every input is a thirteen-crossing literal PD. Join edge labels
through the overpasses, slots one and three with slots numbered
zero through three. Order the resulting components lexicographically
by their increasing label lists; these are arcs $0,\ldots,12$.
For each crossing the input specifies an over-arc $o$ and its two
under-arcs $a,b$.

For each arc the retained input gives a double transposition
$p_i\in A_6$ and a $6\times2$ matrix $Z_i$ over $\F_3$.
Put $P_pe_j=e_{p(j)}$ and
\begin{equation}\label{eq:block}
 M_i=\begin{pmatrix}P_{p_i}&Z_i\\0&I_2\end{pmatrix},
 \qquad V=\F_3^6\oplus\F_3^2.
\end{equation}
All thirteen $8\times8$ matrices are retained literally as well,
so reconstruction from $p_i,Z_i$ is checked independently of that
redundant representation of the input.

\begin{theorem}\label{thm:sharp}
Every $K\in\mathcal K$ has $\mr(K)=b(K)=4$.
Its matrices \eqref{eq:block} define a finite image $H_K$ with
projection $A_6$, a one-dimensional common fixed space and meridian
displacement rank two. The generating rank of its meridian
conjugacy class is exactly four.
\end{theorem}
\begin{proof}
The exact checks establish, for every arc and every crossing,
\[
 M_i^2=I_8,\qquad \rk(M_i-I_8)=2,\qquad M_oM_aM_o=M_b.
\]
The last equation verifies either signed Wirtinger conjugation
convention because all the matrices are involutions. Thus they
give a homomorphism from the literal knot group. Its image is finite
because the field is finite. The four seed permutations generate
$360$ elements. They are even permutations, so their projection is
all of $A_6$.

The checker reconstructs every crossing and verifies the retained
propagation trace from four distinct seed arcs. All thirteen arcs
are reached. Hence the four seed matrices generate the entire image.
Stacking their displacement matrices gives rank seven. The vector
\[
                       c=(1,1,1,1,1,1,0,0)
\]
is fixed by every block matrix, so the common fixed space is
exactly $L=\F_3c$. Lemma~\ref{lem:fixed} gives
\[
                  \mr(K)\ge\left\lceil\frac{8-1}{2}\right\rceil=4.
\]
All meridian images lie in one conjugacy class of $H_K$.
Every generating set from that class has at least four elements,
while the four seed images attain four.

The same four-seed trace is a valid diagrammatic Wirtinger coloring
trace. Blair--Kjuchukova--Velazquez--Villanueva's bridge--Wirtinger
theorem \cite[Definition~1.2 and Theorem~1.3]{bkvv} gives $b(K)\le4$.
Together with $\mr(K)\le b(K)$, this proves the equality.
The exact transfer from each literal diagram to its named census
diagram is a separate input checked as described below.
\end{proof}

For the first case $13a_{304}$, the seed arcs are $0,1,5,9$.
The seed permutations, on $\{0,\ldots,5\}$, are
\[
 p_0=p_1=(2\,3)(4\,5),\qquad
 p_5=(1\,2)(4\,5),\qquad p_9=(0\,2)(3\,4),
\]
and the two nonzero seed blocks are
\[
 Z_0=\begin{pmatrix}0&0\\0&0\\1&2\\2&1\\1&0\\2&0\end{pmatrix},
 \qquad
 Z_1=\begin{pmatrix}0&0\\0&0\\0&0\\0&0\\0&1\\0&2\end{pmatrix},
 \qquad Z_5=Z_9=0.
\]
The remaining nine blocks follow the retained propagation trace.
The unused crossings are checked separately; propagation by itself
would not establish a representation. Row reduction of the four
seed displacement matrices gives rank seven over $\F_3$.

\section{The seven-dimensional images and ordinary rank}

\begin{lemma}\label{lem:quotient}
The action on $\bar V=V/L$ has zero common fixed space and meridian
displacement rank two. In coordinates
\[
 D(w,s)=(w_0-w_5,\ldots,w_4-w_5,s_0,s_1),
\]
its matrix for \eqref{eq:block} is
\[
 \bar M(p,Z)=
 \begin{pmatrix}
  (\mathbf1_{p(j)=i}-\mathbf1_{p(j)=5})_{0\le i,j<5}
    &(Z_{i\alpha}-Z_{5\alpha})_{0\le i<5,\,0\le\alpha<2}\\
 0&I_2
 \end{pmatrix}.
\]
Its meridian conjugacy class also has generating rank exactly four.
\end{lemma}
\begin{proof}
An involution in characteristic three is diagonalizable with
eigenvalues $1,-1$. The line $L$ is fixed, whereas its displacement
image is its $-1$ eigenspace. Their intersection is zero, so
quotienting by $L$ preserves displacement rank.
If $v+L$ is fixed, then $(M_i-I)v\in L$ for every generating
involution. The same zero-intersection argument gives
$(M_i-I)v=0$ for all $i$. Thus $v\in V^{H_K}=L$ and $v+L=0$.
The matrix formula follows by choosing the section $w_5=0$;
the checker also verifies $DM_i=\bar M_iD$ directly.
Lemma~\ref{lem:fixed} and the four projected seed images now give
the class-generating rank.
\end{proof}

\begin{theorem}\label{thm:groups}
For every $K\in\mathcal K$, the eight-dimensional image $H_K$ and
seven-dimensional image $\bar H_K$ have orders
\[
 |H_K|=360\cdot3^6=262440,\qquad
 |\bar H_K|=360\cdot3^4=29160.
\]
Both have ordinary generating rank two. There are exact sequences
\[
 1\longrightarrow C_3^6\longrightarrow H_K\longrightarrow A_6\longrightarrow1,
 \qquad
 1\longrightarrow C_3^4\longrightarrow\bar H_K\longrightarrow A_6\longrightarrow1.
\]
No splitting or mutual isomorphism assertion is needed.
\end{theorem}
\begin{proof}
An element above the identity permutation has an additive
$6\times2$ translation block, so the projection kernel is an
elementary abelian $3$-group. Choose a lifted representative $t_p$
of every projected permutation $p$ by a finite spanning-tree search.
For each $p$ and seed generator $x$, compute
$t_px\,t_{p\pi(x)}^{-1}$. These kernel elements generate the entire
kernel: insert the representatives between consecutive letters of
any word of projected product one and telescope. Inverse edges give
inverses of the corresponding positive-edge elements.
Therefore the kernel is exactly the $\F_3$ span of these translation
blocks. Exact row reduction gives dimension six in every case.

Applying $D$ to the translation blocks gives a kernel span of
dimension four. The permutation action on $\F_3^6/L$ is faithful:
if $P_p-I$ had image in the constant line, each $e_{p(j)}-e_j$
would be constant; its support has at most two coordinates, so it
must vanish. Thus $p=1$. Both quotient projections and order
formulas follow.

Number the four seed images $x_0,x_1,x_2,x_3$ in input order.
For twenty-two cases the pair $x_0x_2,x_1x_3$ generates $H_K$.
For $13a_{629},13a_{3375}$ use $x_0x_1x_2,x_0x_2x_3$ instead.
The same complete Schreier calculation for each pair gives
projection order $360$ and kernel dimension six. Its subgroup
therefore has the full image order and equals $H_K$.
The pair projects to generators of $\bar H_K$ as well.
Their nonabelian quotient $A_6$ excludes ordinary rank one.
\end{proof}

\section{Inputs, independent controls and scope}

The retained package is \path{a6_affine_inputs.json},
\path{verify_a6_affine_certificates.py} and
\path{a6_affine_verification.json}. It contains only the twenty-four
positive certificates. The standard-library checker reconstructs
all $312$ block matrices, checks every involution, displacement rank
and crossing, checks every four-seed trace, and computes the finite
projection groups, Schreier kernels, ordinary generating pairs and
quotient intertwiners. It rejects a changed matrix, a bad seed move,
a corrupted census map and the zero-translation lift.

Each named binding retains a census PD and DT record from
\texttt{snappy\_15\_knots}~1.2.1, read through
Spherogram~2.4.1, with the census database's exact SHA256 and byte
length. The default checker checks the saved crossing-preserving
dart bijection, including rotation, edge gluing and over/under data,
and checks planarity and the one-component condition.
The optional \texttt{--census} mode additionally rereads all named
records and hashes the actual packaged database.
The trusted named-census provenance and the published bridge theorem
are explicit imports, not a numerical knot identification.

The optional \texttt{--cover-homology} control independently builds
the six-sheeted chain complex with all $78$ directed lifted generator
edges. After contracting five tree edges, all $78$ lifted crossing
relations and all $52$ meridian-cycle fillings give a $130\times73$
matrix of rank $71$ over $\F_3$. Thus every retained action has
$\dim H_1(Y;\F_3)=2$. Fixed sheets are filled as well as moving
cycles. A backward lifted generator edge differs from its forward
replacement by the corresponding cycle-filling row. Consequently
the two crossing-conjugation sign conventions give the same filled
row span. This agrees with the numerical bound
$\lceil(6-1+2)/2\rceil=4$, but is not a premise of the direct
matrix proof.

The block matrices were discovered from parabolic cocycles.
The general cocycle-to-affine construction and fixed-space argument
are not asserted new, and the affine realization does not improve
the corresponding homological lower bound numerically.
The explicit full finite images and their distinction between
ordinary generation and generation from the meridian class are
additional checked information. Worldwide priority, membership
outside every known MRC family, formal verification and external
specialist review are not claimed here.

\begin{thebibliography}{9}
\bibitem{bkvv}
R. Blair, A. Kjuchukova, R. Velazquez and P. Villanueva,
Wirtinger systems of generators of knot groups,
Commun. Anal. Geom. 28 (2020), 243--262,
Definition~1.2 and Theorem~1.3, printed p.~246.
The exact coloring move and theorem are imported from the retained
published primary PDF; its full proof is not independently reproved.
\url{https://archive.intlpress.com/site/pub/files/_fulltext/journals/cag/2020/0028/0002/CAG-2020-0028-0002-a002.pdf}.
\end{thebibliography}
\end{document}
